\documentclass[11pt,a4paper]{article}
\usepackage[margin=1in]{geometry}
\usepackage{hyperref}
\usepackage{enumitem}
\usepackage{graphicx}
\usepackage{xcolor}

% -----------------------------------------------------
% bvraghav-tiet-ucs503p-article.sty
% -----------------------------------------------------

% Variable: Lab Instructor
\makeatletter \newcommand{\BvrSetLabInstructor}[1]{%
  \gdef\Bvr@LabInstructor{#1}}
% default
\newcommand{\Bvr@LabInstructor}{Dr. Raghav %
  B. Venkataramaiyer}
% public accessor
\newcommand{\BvrLabInstructorValue}{\Bvr@LabInstructor}
\makeatother

% Add subtitle
\usepackage{titling}
% -- Set title bold
\pretitle{\begin{center}\bfseries\fontsize{18bp}{18bp}\selectfont}
\makeatletter
\newcommand{\BvrSubtitle}[1]{%
  \posttitle{%
    \par\end{center}
    \begin{center}\large#1\end{center}
    \vskip 0.5em}%
}
\makeatother

% Hints in the section title(s)
\newcommand{\BvrHint}[1]{%
  {\normalfont\normalsize\itshape\hfill #1}}

% -----------------------------------------------------

\title{ResearchAgent: An Autonomous AI Agent for Stock and Company Research}

\BvrSetLabInstructor{Paramveer Ma'am}

\BvrSubtitle{%
  Project Proposal for UCS503P  \\
  Submitted to: \BvrLabInstructorValue \\ \bfseries
  Thapar Institute of Engineering and Technology% 
}

\author{
  \begin{tabular}[t]{ccc}
    Ananya Khanduja, CSED & Garv Sharma, CSED & Vishwa Choudhary, CSED \\
    \small \texttt{1024030693} & \small \texttt{1024030695} & \small \texttt{1024030668}
  \end{tabular}
}

\date{\today}

\begin{document}
\maketitle

\section{Higher-order goal}
This project aims to improve access to financial
information and lower the barrier to informed investing
by automating the labour-intensive process of company
research. While the broader motivation is financial
literacy and democratized access to research tooling,
the immediate engineering objective is to translate
\emph{autonomous, tool-using AI agents} into a
measurable, deployable software solution that produces a
structured investment brief for any listed company in
under a minute.

\section{Time-to-value}
\begin{itemize}[leftmargin=0.7cm]
    \item We will deliver meaningful increments quickly by prototyping the core agent loop (a single ticker, a single data source) and validating end-to-end report generation within a short iteration window.
    \item We will prioritize fast feedback over perfection by using weekly iterations and zero-config deployment targets (Render, for both the Go API service and the static React frontend) to reduce release friction.
    \item Success is defined by what can be delivered and measured in the first iteration -- a working plan-act-synthesize loop for one company -- then improved based on report-quality feedback.
\end{itemize}

\section{Problem Statement}
Individual investors and students researching a company or stock face a slow, fragmented process that discourages diligence. Common pain points include:
\begin{itemize}[leftmargin=0.7cm]
    \item \textbf{Fragmentation:} financials, news, and qualitative context live on separate sites (screeners, exchanges, news portals), forcing manual cross-referencing.
    \item \textbf{Cost barriers:} most automated research tools rely on paid financial-data or LLM APIs, which are inaccessible to students and retail investors.
    \item \textbf{Inconsistent structure:} self-compiled notes rarely follow a consistent, comparable format across companies.
    \item \textbf{Stale or incomplete data:} manually gathered snapshots quickly go out of date and often omit recent news or risk factors.
\end{itemize}

\textbf{Why this matters now (contextual relevance):} With free-tier LLMs and public data sources now capable of near-real-time synthesis, a zero-paid-API research agent can meaningfully close the access gap for individual users without institutional data subscriptions.

\section{Proposed Solution}
\subsection{Overview}
Build a \textbf{web-based} autonomous research agent, ``ResearchAgent'', with the following components:
\begin{itemize}[leftmargin=0.7cm]
    \item \textbf{Company/ticker input portal:} user submits a company name or stock ticker (e.g. \texttt{TCS}, \texttt{BLS}).
    \item \textbf{Automated ticker resolution:} agent infers the correct NSE/BSE ticker symbol without manual lookup.
    \item \textbf{Autonomous data gathering:} agent pulls live financial data, web search results, and recent news for the target company.
    \item \textbf{Structured report synthesis:} agent compiles all gathered context into a consistent, six-section investment brief.
    \item \textbf{Report viewer:} a dedicated page renders the completed brief per ticker for review and sharing.
\end{itemize}

\subsection{Core workflow}
\begin{enumerate}[leftmargin=0.7cm]
    \item User submits a company name and the agent resolves it to a ticker symbol.
    \item Agent plans a sequence of research steps (search, financials, news) tailored to that company.
    \item Agent dispatches tool calls, observes results, and retries with a different query or tool if a source returns insufficient data.
    \item Agent synthesizes all gathered context into the six-section report and renders it to the user.
\end{enumerate}

\subsection{Operational constraints for an educational, zero-budget setting}
\begin{itemize}[leftmargin=0.7cm]
    \item Rely only on free-tier services: a single LLM API key (Gemini) is the sole required credential.
    \item Avoid dependence on paid financial-data providers; use a free-tier financial-data API (Twelve Data) as an optional primary source, with public scrapers and RSS feeds as layered fallbacks.
    \item Design for deployability using zero-config hosting (Render) for both the stateless Go backend and the static React frontend, avoiding any managed database or paid infrastructure.
\end{itemize}

\section{Solution Approach}
This is an engineering course project, so we focus on agent orchestration, tool-integration reliability, and system correctness rather than novel model research.

\subsection{Agent orchestration}
The core agent loop runs entirely inside a single API route and follows a \emph{Plan $\rightarrow$ Act $\rightarrow$ Observe $\rightarrow$ Reflect $\rightarrow$ Synthesise} pattern:
\begin{itemize}[leftmargin=0.7cm]
    \item The LLM first generates a step-by-step research plan for the given company.
    \item The agent dispatches calls to three tools -- search, financials, and news -- based on that plan.
    \item After each tool call, the agent reflects on whether results are sufficient and retries with an alternate query or tool if not, without claiming guaranteed completeness.
    \item Gathered context from all tool calls is passed to the LLM for final synthesis into the report.
\end{itemize}

\subsection{Web architecture}
A decoupled Go backend and React frontend, communicating over a plain HTTP/JSON API:
\begin{itemize}[leftmargin=0.7cm]
    \item \textbf{Frontend:} a responsive UI (React with Vite and TypeScript, Tailwind CSS) for ticker input and report viewing, deployed as a static site independent of the backend.
    \item \textbf{Backend API:} a Go HTTP server, built on the standard library alone, exposing a core agent route (\texttt{/api/agent}) plus dedicated tool routes for search, financials, and news.
    \item \textbf{External data sources:} Twelve Data as an optional, documented primary source for financials, with a crumb-authenticated Yahoo Finance quote endpoint and a Screener.in scraper as layered fallbacks; DuckDuckGo HTML search with a Google News RSS fallback for web search; Google News RSS for dedicated news retrieval.
\end{itemize}

\subsection{Search and data resilience}
Because the system depends on unofficial/public sources, resilience is built in at the tool level:
\begin{itemize}[leftmargin=0.7cm]
    \item Web search attempts DuckDuckGo first, falls back to Google News RSS if rate-limited, and returns an empty result gracefully so the agent can retry rather than fail.
    \item Financial data resolution attempts Twelve Data first when configured (a single documented API covering both global and NSE/BSE-listed tickers), then falls back through a crumb-authenticated Yahoo Finance quote (with \texttt{.NS}/\texttt{.BO} suffix logic for Indian exchanges), an unauthenticated Yahoo chart endpoint, and finally a Screener.in scrape, returning whichever source succeeds first as a single financial-data object.
\end{itemize}

\subsection{CI/CD and fast delivery enabler}
CI/CD will be used to:
\begin{itemize}[leftmargin=0.7cm]
    \item Automatically build, vet, and test both the Go backend (\texttt{go build}, \texttt{go vet}, unit tests) and the React frontend (\texttt{tsc}, \texttt{vite build}) on every change.
    \item Produce repeatable deployment artifacts via Render's zero-config pipelines for a Go web service and a static frontend site.
    \item Enable safe and frequent releases of incremental features (new report sections, new data sources) across both services independently.
\end{itemize}

\section{Evaluation Criterion}
We define success using measurable metrics tied to the core research pipeline.

\subsection{Primary evaluation metric}
\textbf{Time-to-report (TTR):} time between a user submitting a company name and receiving a complete, six-section report.
\begin{itemize}[leftmargin=0.7cm]
    \item Target: median TTR $\le$ 60 seconds for a resolved NSE/BSE ticker.
    \item Attribution: measured server-side from request receipt to synthesis completion in the \texttt{/api/agent} route.
\end{itemize}

\subsection{Secondary evaluation metrics}
\begin{itemize}[leftmargin=0.7cm]
    \item \textbf{Ticker resolution accuracy:} percentage of company-name inputs correctly resolved to the right NSE/BSE ticker without manual correction.
    \item \textbf{Data completeness:} percentage of reports where all six sections are populated with non-empty, source-grounded content (no section left blank due to tool failure).
    \item \textbf{Fallback trigger rate:} frequency with which the search/news fallback path is invoked, as an indicator of primary-source reliability.
    \item \textbf{Reliability:} availability of the deployed agent endpoint during testing windows (e.g. $\ge$ 99\% uptime on Render free tier).
\end{itemize}

\subsection{Pilot validation plan}
\begin{itemize}[leftmargin=0.7cm]
    \item Run the agent against a curated list of 20--30 companies spanning large-cap, mid-cap, and small-cap NSE/BSE listings.
    \item Manually spot-check a sample of generated reports against source data (Screener.in, recent news) to validate factual grounding before broader use.
\end{itemize}

\section{Scalability}
The proposition should scale in both theoretical and practical senses.

\begin{enumerate}[leftmargin=0.7cm]
    \item \textbf{Scalable (theoretically):} The system should support growth in concurrent analyses and request volume by:
    \begin{itemize}[leftmargin=0.7cm]
        \item decoupling the frontend from the stateless agent API route,
        \item batching and caching tool responses (financials, news) per ticker to reduce redundant calls,
        \item keeping each tool call independently retryable and timeout-bounded.
    \end{itemize}
    
    \item \textbf{Operationally deployable (practically):} The system should run entirely on free-tier infrastructure:
    \begin{itemize}[leftmargin=0.7cm]
        \item a stateless Go API service and a static React frontend, both hosted on Render's free tier with no managed database to provision,
        \item a single required LLM credential (Gemini, with OpenRouter as an optional backup path) and optional free-tier keys (Twelve Data, Exa) that gracefully degrade to public scrapers when absent,
        \item CI/CD pipeline for automatic staging/production deployment of both services on every push.
    \end{itemize}
\end{enumerate}

\section{Engine availability heuristic}
A mature set of ``engines'' (tooling/services) is available to implement this as a web-based project:
\begin{itemize}[leftmargin=0.7cm]
    \item Go, using only its standard library, for the backend HTTP API; React with Vite and TypeScript for the frontend.
    \item Tailwind CSS v4 for a dark-mode-first responsive frontend.
    \item Google Gemini 2.0 Flash-Lite as the primary LLM, with an OpenRouter-routed Gemini 2.0 Flash wrapper as an optional fallback path.
    \item Twelve Data as an optional, documented free-tier financial-data API, with a crumb-authenticated Yahoo Finance quote endpoint and a Screener.in scraper as layered fallbacks, none requiring a paid API key.
    \item DuckDuckGo HTML search and Google News RSS for web search and news, with layered fallback, parsed using Go's standard library.
    \item Render free tier for zero-config hosting and deployment of both the backend service and the static frontend.
\end{itemize}
We will use these mature components to focus on agent orchestration, data-source reliability, and report quality rather than inventing new infrastructure.

\section{Project scope and deliverables}
\subsection{Initial deliverable}
\begin{itemize}[leftmargin=0.7cm]
    \item Company/ticker input form and automated ticker resolution
    \item Core agent loop: plan generation, single-pass tool dispatch, and synthesis for one company at a time
    \item Financials tool (Twelve Data as optional primary source, with Yahoo Finance and Screener.in as layered fallbacks) and basic report rendering
    \item Basic logging and timestamps for time-to-report measurement
    \item CI: build + route-level tests + linting
    \item CD to Render staging with a single push-to-deploy pipeline
\end{itemize}

\subsection{Subsequent deliverables}
\begin{itemize}[leftmargin=0.7cm]
    \item Search and news tools with layered fallback (DuckDuckGo, Google News RSS)
    \item Reflect/retry logic so the agent re-queries when tool results are insufficient
    \item Full six-section report: overview, financial snapshot, news, competitive landscape, risk factors, and bull/bear investment summary
    \item Report history/viewer page per ticker and lightweight caching for repeat queries
\end{itemize}

\section{Risks and mitigations}
\begin{itemize}[leftmargin=0.7cm]
    \item \textbf{Unofficial data source instability:} Yahoo Finance's quote endpoint and Screener.in are unofficial/scraped sources that can change or be blocked without notice; mitigate by preferring Twelve Data's documented API when configured, and by layering an unauthenticated Yahoo chart fallback and graceful degradation rather than hard failures on the remaining unofficial paths.
    \item \textbf{Rate limiting:} DuckDuckGo scraping may be throttled under load; mitigate with the Google News RSS fallback and request pacing.
    \item \textbf{Hallucinated or unverifiable claims:} LLM synthesis may overstate confidence; mitigate by grounding synthesis strictly in retrieved tool output and clearly labelling the report as informational, not financial advice.
    \item \textbf{Free-tier quota limits:} Gemini/OpenRouter and Twelve Data free tiers cap request volume (e.g. Twelve Data: 800 requests/day, 8/minute); mitigate by using the OpenRouter fallback path, treating Twelve Data as an optional source with automatic fallback, and caching recent reports.
\end{itemize}

\section{Summary}
This project proposes a deployable, zero-paid-API platform -- a Go backend paired with a React frontend -- that autonomously researches a company or stock ticker and produces a structured investment brief. It meets the course criteria by emphasizing time-to-value through rapid, tool-by-tool iteration, implementing CI/CD to support frequent safe releases across both services, and using measurable evaluation criteria (especially time-to-report and data completeness). It remains focused on engineering agent orchestration and data-source reliability rather than research-driven model novelty.

\end{document}
